\documentclass{article}
\usepackage{mathtools} 
\usepackage{fontspec}
\usepackage[UTF8]{ctex}
\usepackage{amsthm}
\usepackage{mdframed}
\usepackage{xcolor}
\usepackage{amssymb}
\usepackage{amsmath}


% 定义新的带灰色背景的说明环境 zremark
\newmdtheoremenv[
  backgroundcolor=gray!10,
  % 边框与背景一致，边框线会消失
  linecolor=gray!10
]{zremark}{说明}

% 通用矩阵命令: \flexmatrix{矩阵名}{元素符号}{行数}{列数}
\newcommand{\flexmatrix}[4]{
  \[
  #1 = \begin{pmatrix}
    #2_{11}     & #2_{12}     & \cdots & #2_{1#4}   \\
    #2_{21}     & #2_{22}     & \cdots & #2_{2#4}   \\
    \vdots      & \vdots      & \ddots & \vdots     \\
    #2_{#31}    & #2_{#32}    & \cdots & #2_{#3#4}
  \end{pmatrix}
  \]
}

% 简化版命令（默认矩阵名为A，元素符号为a）: \quickmatrix{行数}{列数}
\newcommand{\quickmatrix}[2]{\flexmatrix{A}{a}{#1}{#2}}

\begin{document}
\title{3.1}
\author{张志聪}
\maketitle

\begin{zremark}
  设$\{A_k\}_{k = 1}^{\infty} \in \mathbb{C}^{n \times n}$中的矩阵序列，
  即$\lim\limits_{k \to \infty} A_k = A$，
  则
  \begin{itemize}
    \item (1) $\lim\limits_{k \to \infty} det(A_k) = det(A)$;
    \item (2) $\lim\limits_{k \to \infty} A_k^{*} = A^{*}$，其中$A_k^{*}, A^{*}$是
          $A_k, A$的伴随矩阵.
  \end{itemize}
\end{zremark}

\textbf{证明：}
\begin{itemize}
  \item (1) 设$A_k = (a_{ij}^{(k)}), A = (a_{ij})$，由$A_k \to A$可得
        \begin{align*}
          \lim\limits_{k \to \infty} a_{ij}^{(k)} = a_{ij}
        \end{align*}

        通过对$n$进行归纳，我们来完成证明。

        $n = 1$时，$det(A_k) = a_{11}^{(k)}$结论是显然的;

        归纳假设$n - 1$时，命题成立;

        现在需要证明$n$时，命题成立。对$A_k$第一行进行展开，我们有
        \begin{align*}
          det(A_k)
           & = \sum\limits_{j = 1}^n a_{1j}^{(k)} C_{1j}^{(k)}
        \end{align*}
        其中$C_{1j}^{(k)} = (-1)^{1 + j} M_{1j}^{(k)}$，
        $M_{1j}^{(k)}$是$n - 1$阶方阵的余子式。

        由$A_k \to A$，我们有
        \begin{align*}
          a_{1j}^{(k)} \to a_{1j}
        \end{align*}
        由归纳假设可知，
        \begin{align*}
          M_{1j}^{(k)} \to M_{1j}
        \end{align*}
        故
        \begin{align*}
          \sum\limits_{j = 1}^n a_{1j}^{(k)} C_{1j}^{(k)} = det(A)
        \end{align*}
        归纳完成，命题得证。
  \item (2) 可以类似的归纳证明。
\end{itemize}

\begin{zremark}
  例3的补充
  \begin{align*}
    \left\| (I - A)^{-1} - \sum\limits_{k=0}^m A^k \right\|
     & = \left\| \sum\limits_{k=0}^\infty A^k - \sum\limits_{k=0}^m A^k \right\| \\
     & = \left\| \sum\limits_{k=m + 1}^\infty A^k\right\|                        \\
     & \leq \sum\limits_{k=m + 1}^\infty \| A \|^k                               \\
     & = \frac{\|A\|^{m + 1}}{1 - \|A\|}
  \end{align*}
  注意：$\|A\| < 1$才成立。
\end{zremark}

\begin{zremark}
  证明：

  (I) $\forall A \in \mathbb{C}^{n \times n}$，总有
  \begin{itemize}
    \item (1) $sin (-A) = -sin A, cos(-A) = cos A$;
  \end{itemize}
\end{zremark}

\textbf{证明：}
\begin{align*}
  - sin A
   & = - \sum\limits_{k = 0}^\infty \frac{(-1)^k A^{2k + 1}}{(2k + 1)!}   \\
   & = \sum\limits_{k = 0}^\infty \frac{(-1)^{k+1} A^{2k + 1}}{(2k + 1)!} \\
\end{align*}
又
\begin{align*}
  sin (-A)
   & = \sum\limits_{k = 0}^\infty \frac{(-1)^k (-A)^{2k + 1}}{(2k + 1)!}           \\
   & = \sum\limits_{k = 0}^\infty \frac{(-1)^k (-1)^{2k + 1}A^{2k + 1}}{(2k + 1)!} \\
   & = \sum\limits_{k = 0}^\infty \frac{(-1)^{k + 1}A^{2k + 1}}{(2k + 1)!}         
\end{align*}
故
\begin{align*}
  -sin A = sin (-A)
\end{align*}

\end{document}